% 从旧数学附录迁出的计数、逆元与同余说明；不新增算法。
\section{矩形 GCD 计数：缩放后容斥}\label{knowledge-coprime-rectangles}\index{coprime pairs}\index{矩形 GCD 计数}
前缀矩形互质数对数
\[
F(A,B)=\sum_{i=1}^{\min(A,B)}\mu(i)\left\lfloor\frac Ai\right\rfloor
\left\lfloor\frac Bi\right\rfloor,
\]
按两个商同时分块，右端点 $r=\min(A/\lfloor A/l\rfloor,B/\lfloor B/l\rfloor)$，用 $\mu$ 前缀和；
只需筛到 $\min(A,B)$。$A=0$ 或 $B=0$ 时为 0。
矩形 $[a,b]\times[c,d]$ 内 $\gcd(x,y)=k$ 的对数为
\[
\begin{aligned}
&F(\lfloor b/k\rfloor,\lfloor d/k\rfloor)
-F(\lfloor(a-1)/k\rfloor,\lfloor d/k\rfloor)\\
&-F(\lfloor b/k\rfloor,\lfloor(c-1)/k\rfloor)
+F(\lfloor(a-1)/k\rfloor,\lfloor(c-1)/k\rfloor),
\end{aligned}
\]
必须先减 1 再除以 $k$。答案不超过面积，可存有符号 64 位；分块累加用 128 位。

\section{幂的约数和：不除以 \texorpdfstring{$p-1$}{p-1}}\label{knowledge-divisor-power-sum}\index{divisor sum}\index{geometric series}
若 $a=\prod p_i^{e_i}$，则 $\sigma(a^b)=\prod_i(1+p_i+\cdots+p_i^{e_i b})$。
合数模数或 $p\equiv1$ 时分母不可逆，不能用几何级数除法公式。
令 $G(a,n)=(a^n,\sum_{i=0}^{n-1}a^i)$，$G(a,0)=(1,0)$；由 $h=\lfloor n/2\rfloor$ 的 $(P,S)$，
偶数项得 $(P^2,S(1+P))$，奇数项得 $(P^2a,S(1+P)+P^2)$。时间 $O(\log(n+1))$，任意正模数。
参数是项数，约数和须传 $e_ib+1$；$e_ib$ 可能超过 64 位，模乘在 128 位中做。

\section{任意数组的批量逆元}\label{knowledge-batch-inverse}\index{批量逆元}
设 $s_0=1$，$s_i=\prod_{j<i}a_j\pmod m$，只对 $s_n$ 求一次逆元，再从右向左：若 $t=s_{i+1}^{-1}$，则
\[
a_i^{-1}=t s_i\pmod m,\qquad s_i^{-1}=t a_i\pmod m.
\]
先保存原 $a_i$ 再覆盖。总乘积与 $m$ 互素当且仅当每项互素，模数无需为素数。
实现见第~\pageref{compact-batch_inverse}~页：无解返回空的可选值，空输入返回空向量，模数 1 返回全零。

\section{连续整数的线性逆元表}\label{knowledge-linear-inverses}\index{线性逆元表}
素数 $p$、$1<i<p$ 时
\[
i^{-1}\equiv-\lfloor p/i\rfloor\,(p\bmod i)^{-1}\pmod p,
\]
从 $1^{-1}=1$ 递推，$O(N)$。实现见第~\pageref{compact-inverse_table}~页，第 0 项占位，允许 $N=0$；
要求 $N<p$。合数模数不能照搬（$p\bmod i$ 可能不可逆），单个逆元用扩展欧几里得（$\gcd(a,m)=1$ 才有解）
或素数下的 $a^{p-2}$。

\section{线性同余：最小解与最小周期}\label{knowledge-linear-congruence}\index{linear congruence}\index{线性同余方程}
实现见第~\pageref{compact-linear_congruence}~页，解 $ax+b\equiv0\pmod m$。
$a$ 归一化后令 $g=\gcd(a,m)$；有解当且仅当 $g\mid b$，若 $au+mv=g$，则 $x_0=u(-b/g)$，
全部解为 $x_0+(m/g)t$，返回 $[0,m/g)$ 内最小解与最小正周期 $m/g$。
$a=0$ 且 $m\mid b$ 或 $m=1$ 时返回 $(0,1)$。$b$ 可为 \texttt{LLONG\_MIN}（内部用 \texttt{-i128(b)}）；
题目为 $ax\equiv c$ 时先把 $c$ 归一化再取负传入。exgcd 见第~\pageref{knowledge-diophantine-bounds}~页。
追及问题 $(m-n)t+(x-y)\equiv0\pmod L$ 直接传速度差与起点差，不取绝对值。

\section{扩展 CRT：逐式合并}\label{knowledge-excrt}\index{exCRT}\index{扩展中国剩余定理}
实现见第~\pageref{compact-crt_merge}~页的 \texttt{crt\_merge}：从 $r=0,m=1$ 起依次合并 $x\equiv b\pmod n$，
一旦无解即停止。设 $g=\gcd(m,n)$，$mu+nv=g$；若 $g\nmid(b-r)$ 无解，否则
\[
q=n/g,\qquad t\equiv\frac{b-r}{g}u\pmod q,\qquad 0\le t<q,
\]
新状态为 $((r+mt)\bmod mq,\ mq)$。模数不必互素。
要求正模数且最终周期不超过有符号 64 位；中间量用 128 位。无解返回假，周期越界抛异常。
时间 $O(k\log M)$，递归栈 $O(\log M)$。

\section{Garner：总周期可以超过整数范围}\label{knowledge-garner}\index{Garner}\index{混合进制}
模数 $m_i$ 两两互素，$P=\prod m_i$，把解写成混合进制
\[
x=d_0+d_1m_0+d_2m_0m_1+\cdots,\qquad 0\le d_i<m_i.
\]
对每个后续模数维护已定部分余数 $v_j$ 与位权 $c_j=\prod_{t<i}m_t\bmod m_j$，
$d_i=(b_i-v_i)c_i^{-1}\bmod m_i$，然后 $v_j\gets v_j+c_jd_i$，$c_j\gets c_jm_i$。
所得 $x\in[0,P)$ 为最小非负解；模数 1 可直接参与。实现见第~\pageref{compact-garner}~页。
求 $x\bmod T$ 时从高位 Horner，$T$ 无需与 $m_i$ 互素；结果只是规范解的余数，
多模 NTT 还原整数时须保证真实值在 $[0,P)$ 内。
